% GENERATED FILE. Edit canonical lessons, then run npm run generate:books.
% canonical-source-hash: fnv1a64:d10e9ba4f21d9bce
% canonical-lessons: FR-C155-repos, FR-C155-sommeil, FR-C155-reve, FR-C155-angoisse, FR-C155-amour

\chapter{Rest, Sleep, Dream, Anxiety, Love}
\label{ch:fr-rest-sleep-dream-anxiety-love}
\hlchaptermodality{155}

\begin{chapteropening}
\textbf{By the end of this chapter:} \emph{I can say rest, sleep, a dream, anxiety and love.}

Complete the last lesson of chapter 155: I can say rest, sleep, a dream, anxiety and love.
\end{chapteropening}

\section[le repos]{le repos — rest}

\label{lesson:FR-C155-repos}



\begin{quote}
Before the new one: say the French for the liver, then the French for a kidney.
\end{quote}

\subsection*{You'll want to know: le repos}
\textbf{le repos} — \textquotedblleft{}rest\textquotedblright{}.

\begin{cousinweb}
\textbf{repos} — formed from the verb for \textquotedblleft{}to rest\textquotedblright{}, $\leftarrow$ Late Latin \emph{repausare}, \textquotedblleft{}to rest\textquotedblright{}, from \emph{pausare}, \textquotedblleft{}to pause\textquotedblright{}. (EtymDB 2.1)
\end{cousinweb}

\begin{grammarlens}[title={What you've built}]
One more word for this chapter.
\end{grammarlens}

\subsection*{Your turn}
\begin{itemize}
\raggedright
  \item \emph{Say it:} \emph{le repos}
  \item \emph{Say it:} \emph{le repos}, once more
  \item \emph{Say it:} say \emph{le repos}
  \item \emph{Recall:} say the French for the liver, then the French for a kidney, then say \emph{le repos} again
\end{itemize}

\begin{tcolorbox}[breakable,colback=teal!4,colframe=teal!35!black,title={Before you move on}]
What does le repos mean? (\textquotedblleft{}Rest\textquotedblright{}.) Say it once more.
\end{tcolorbox}
\section[le sommeil]{le sommeil — sleep}

\label{lesson:FR-C155-sommeil}



\begin{quote}
Before the new one: say the French for a kidney, then the French for rest.
\end{quote}

\subsection*{You'll want to know: le sommeil}
\textbf{le sommeil} — \textquotedblleft{}sleep\textquotedblright{}.

\begin{cousinweb}
\textbf{sommeil} $\leftarrow$ Vulgar Latin \emph{*somniculus}, \textquotedblleft{}a little sleep\textquotedblright{}, from Latin \emph{somnus}, \textquotedblleft{}sleep\textquotedblright{} (Lewis \& Short) — inherited.
\end{cousinweb}

\begin{grammarlens}[title={What you've built}]
One more word for this chapter.
\end{grammarlens}

\subsection*{Your turn}
\begin{itemize}
\raggedright
  \item \emph{Say it:} \emph{le sommeil}
  \item \emph{Say it:} \emph{le sommeil}, once more
  \item \emph{Say it:} say \emph{le sommeil}
  \item \emph{Recall:} say the French for a kidney, then the French for rest, then say \emph{le sommeil} again
\end{itemize}

\begin{tcolorbox}[breakable,colback=teal!4,colframe=teal!35!black,title={Before you move on}]
What does le sommeil mean? (\textquotedblleft{}Sleep\textquotedblright{}.) Say it once more.
\end{tcolorbox}
\section[le rêve]{le rêve — a dream}

\label{lesson:FR-C155-reve}



\begin{quote}
Before the new one: say the French for rest, then the French for sleep.
\end{quote}

\subsection*{You'll want to know: le rêve}
\textbf{le rêve} — \textquotedblleft{}a dream\textquotedblright{}.

\begin{cousinweb}
\textbf{rêve} — formed from \emph{rêver}, \textquotedblleft{}to dream\textquotedblright{}.
\end{cousinweb}

\begin{grammarlens}[title={What you've built}]
One more word for this chapter.
\end{grammarlens}

\subsection*{Your turn}
\begin{itemize}
\raggedright
  \item \emph{Say it:} \emph{le rêve}
  \item \emph{Say it:} \emph{le rêve}, once more
  \item \emph{Say it:} say \emph{le rêve}
  \item \emph{Recall:} say the French for rest, then the French for sleep, then say \emph{le rêve} again
\end{itemize}

\begin{tcolorbox}[breakable,colback=teal!4,colframe=teal!35!black,title={Before you move on}]
What does le rêve mean? (\textquotedblleft{}A dream\textquotedblright{}.) Say it once more.
\end{tcolorbox}
\section[l'angoisse]{l'angoisse — anxiety}

\label{lesson:FR-C155-angoisse}



\begin{quote}
Before the new one: say the French for sleep, then the French for a dream.
\end{quote}

\subsection*{You'll want to know: l'angoisse}
\textbf{l'angoisse} — \textquotedblleft{}anxiety\textquotedblright{}.

\begin{cousinweb}
\textbf{angoisse} $\leftarrow$ Latin \emph{angustia}, \textquotedblleft{}narrowness, distress\textquotedblright{}, from \emph{angustus}, \textquotedblleft{}narrow\textquotedblright{} — inherited. (EtymDB 2.1)
\end{cousinweb}

\begin{grammarlens}[title={What you've built}]
One more word for this chapter.
\end{grammarlens}

\subsection*{Your turn}
\begin{itemize}
\raggedright
  \item \emph{Say it:} \emph{l'angoisse}
  \item \emph{Say it:} \emph{l'angoisse}, once more
  \item \emph{Say it:} say \emph{l'angoisse}
  \item \emph{Recall:} say the French for sleep, then the French for a dream, then say \emph{l'angoisse} again
\end{itemize}

\begin{tcolorbox}[breakable,colback=teal!4,colframe=teal!35!black,title={Before you move on}]
What does l'angoisse mean? (\textquotedblleft{}Anxiety\textquotedblright{}.) Say it once more.
\end{tcolorbox}
\section[l'amour]{l'amour — love}

\label{lesson:FR-C155-amour}



\begin{quote}
Before the new one: say the French for a dream, then the French for anxiety.
\end{quote}

\subsection*{You'll want to know: l'amour}
\textbf{l'amour} — \textquotedblleft{}love\textquotedblright{}.

\begin{cousinweb}
\textbf{amour} $\leftarrow$ Old French \emph{amor} $\leftarrow$ Latin \emph{amor}, \textquotedblleft{}love\textquotedblright{} — inherited (EtymDB 2.1). Its \emph{-our}, where French usually has \emph{-eur} as in \emph{douleur}, is often put down to the love poetry of the south.
\end{cousinweb}

\begin{grammarlens}[title={What you've built}]
One more word for this chapter.
\end{grammarlens}

\subsection*{Your turn}
\begin{itemize}
\raggedright
  \item \emph{Say it:} \emph{l'amour}
  \item \emph{Say it:} \emph{l'amour}, once more
  \item \emph{Say it:} say \emph{l'amour}
  \item \emph{Recall:} say the French for a dream, then the French for anxiety, then say \emph{l'amour} again
\end{itemize}

\begin{tcolorbox}[breakable,colback=teal!4,colframe=teal!35!black,title={Before you move on}]
What does l'amour mean? (\textquotedblleft{}Love\textquotedblright{}.) Say it once more.
\end{tcolorbox}
